
\part{항목별 입력 방법}

\chapter{건물 형상}

\section{평가 범위}
건축물의 에너지절약설계기준에 따른 단열선 내의 냉난방 공간을 평가 대상으로 하며, 평면도상에 명기된 실들을 모델링 범위에 포함한다. (예: 냉난방 설비가 설치되지 않는 미단열 공간인 창고 및 주차장 등은 모델링 범위에서 제외한다.) 

\section{도면의 해석}
각 존의 바닥, 천장 및 외피(벽체, 창호 등) 면적은 해당 부위 구조체의 벽체 중심선을 기준으로 산정하여 입력한다.

\section{건물의 조닝}

\section{단순하지 않은 형상의 처리}

\subsection{곡면 외피의 모델링}

원형 평면은 8방위로 나누어 모델링하되, 각 분할 면에 해당되는 방위각을 입력한다. 곡선면의 경우 외벽의 법선면을 기준으로 각도(0-360, 1°단위)를 입력하며, 도면상 정밀한 각도 산출이 어려운 경우 8개 방위로 간략화하여 각도를 입력한다.

\chapter{설비 시스템}

\section{냉난방 시스템}

용량이 큰 설비를 하나만 입력하세요.

용량이나 성능은 모르면 입력하지 마세요. 기본값은 Technical Reference Manual에 나와 있습니다.

\section{신재생 에너지}

\subsection{효율}

효율은 ...를 보고 입력하는 것이 원칙이다.

관계도서가 확인되지 않을 떄에는, 시스템의 종류에 따라 \cref{tbl:PVpanel-efficiency}의 값으로 입려 가능하다. 여기서, 건물의 벽면에 부착된 BIPV는 밀착형으로, 지면과 떨어져 독립적으로 설치된 PV는 후면통풍형으로 본다.

\begin{table}[htbp]
    \caption{태양광 패널의 효율이 명시되지 않는 경우 입력값}
    \label{tbl:PVpanel-efficiency}
    \centering
    \begin{tabular}{cccc}
        \toprule
        패널의 종류 & \multicolumn{3}{c}{환기방법} \\
        & 밀착형 & 후면통풍형 & 기계환기형 \\
        \midrule
        단결정 & 12.6\% & 13.5\% & 14.4\% \\
        다결정 & 11.2\% & 12.0\% & 12.8\% \\
        \bottomrule
    \end{tabular}
\end{table}

여기서, 태양광 패널의 종류를 확인할 수 없는 경우 단결정으로 본다.